\section{Week 6: region-local mu equivalence}

\subsection{Why whole-memory equality is wrong}

The Sign hash runs before signature publication, whereas the Verify hash runs
after the 1474-byte signature and 8-byte length writes. Thus the two global
memories are generally unequal even when every hash input byte is preserved.
The correct relation is
\[
  \mathsf{region\_eq}(M_1,M_2,b,n)
  \eqdef
  \forall i\in[0,n),\;M_1[b+i]=M_2[b+i].
\]

The Week 6 data flow is:
\[
\begin{array}{c}
 M_{KG}\;\xrightarrow{\text{SK import}}\;\mathsf{SignMu}^{64}\\
 \quad\searrow\;\text{signature publication}\;\to M_{after\,Sign}\\
 \qquad\quad\text{preserved VK/pre/message}\;\to\mathsf{VerifyMu}^{32}.
\end{array}
\]

\subsection{Generated procedure theorem}

Theory \path{RegionLocalMuEquivalence.ec} proves locality of the actual
generated pre/message absorb loops, including the load-index and W64 tail
invariants. Key absorption is separate: Sign reads the local packed SK array,
while Verify reads the external VK bytes.

The final theorem
\begin{quote}
\small
\path{RegionLocalMuTop.sign_verify_generated_raw_mu_prefix_regionwise}
\end{quote}
has the actual generated procedures
\procname{_sf_mu_rawpre} and \procname{__verify_hash_mu} on its two sides. Its
premises expose mode-2 lengths, \texttt{raw\_prelen}, non-wrapping pre/message
regions, the local-SK/external-VK 992-byte prefix, and byte equality only on
the two address ranges read by the hash. Its conclusion is
\[
  \take_{32}(\mu_{\mathrm{Sign}}^{64})
  = \mu_{\mathrm{Verify}}^{32}.
\]

No SHAKE prefix or ``outside writes do not matter'' axiom is used. Init,
Keccak permutation, finalize, and both squeeze procedures are the pinned
generated programs. The generated focused hashes remain
\noindent\path{SignMuHashTarget.ec}:
\begin{verbatim}
c2232e64adf71ee54763f51e11862938c8caa5ca278af30e4c83756091e83aa7
\end{verbatim}
\noindent\path{VerifyMuHashTarget.ec}:
\begin{verbatim}
fcee969af0b567447c6c9b9c53e29245375ec2bfd54ae11afecdf0a6ebc0ebb7
\end{verbatim}

\subsection{Address subtlety}

The predicate \texttt{canonical\_region(a,n)} allows the empty boundary case
\(a=2^{64},n=0\). It therefore cannot by itself justify
\(\operatorname{to\_uint}(\operatorname{of\_int}(a))=a\). The caller bridge
keeps \path{canonical_ui64_address} premises for each address or length
whose exact round trip is used. This prevents a silent modular-wrap premise.

Status: \status{PROVED} for mode-2 raw/internal hashing.
